\section*{Hoofdstuk \ref{ch:b3:nervous-systems} --- Organisatie van zenuwstelsels}

\begin{solution}{exo:b3:nervous-systems:1}
\omterm{def:b3:nervous-systems:cells}{Astrocyten}: bufferen kalium en glutamaat, voeden neuronen, wekken de
bloed--hersenbarrière op, regelen de plaatselijke doorbloeding.
\omterm{def:b3:nervous-systems:cells}{Oligodendrocyten}: myeliniseren axonen in het centrale zenuwstelsel.
Schwann-cellen: myeliniseren perifere axonen (en leiden hun
regeneratie). \omterm{def:b3:nervous-systems:cells}{Microglia}: de \omterm{def:b3:innate-immunity:innate}{macrofagen} van het brein, die synapsen
snoeien en puin opruimen. (Ependymcellen bekleden de ventrikels en
maken het \omterm{prop:b3:nervous-systems:barrier-budget}{hersenvocht}.)
\end{solution}

\begin{solution}{exo:b3:nervous-systems:2}
De kleuring zwartte enkele neuronen volledig en liet de rest
onzichtbaar, zodat één cel in haar geheel zichtbaar werd. Cajal besloot
dat neuronen afzonderlijke cellen zijn die elkaar raken zonder te
versmelten, met een vaste richting van de stroom, van dendrieten naar
axon; Golgi besloot dat zij een doorlopend netwerk vormden. Cajal had
gelijk; de elektronenmicroscoop liet de synaptische spleet in 1954 zien.
\end{solution}

\begin{solution}{exo:b3:nervous-systems:3}
Een tik rekt de quadriceps; zijn spierspoelen vuren; de Ia-afferenten
komen het merg binnen en schakelen rechtstreeks over op de
motorneuronen van de quadriceps, die vuren en de spier samentrekken,
terwijl een \omterm{def:b3:nervous-systems:cells}{interneuron} de motorneuronen van de hamstring remt. Eén
synaps, korte axonen bij \qty{80}{m/s}: ongeveer \qty{30}{ms}. Een
bewuste trap gaat door het brein --- waarneming, beslissing, motorische
planning, afdalende banen, meerdere synapsen --- en duurt
\qtyrange{200}{300}{ms}.
\end{solution}

\begin{solution}{exo:b3:nervous-systems:4}
Sympathisch (noradrenaline bij het doelorgaan): hart sneller en
krachtiger, pupil verwijd, darmbeweging en -afscheiding verminderd,
luchtwegen verwijd. Parasympathisch (acetylcholine): hart vertraagd,
pupil vernauwd, darmbeweging en -afscheiding verhoogd, luchtwegen
vernauwd.
\end{solution}

\begin{solution}{exo:b3:nervous-systems:5}
$\lambda = \sqrt{R_{m}d/4R_{i}} = \sqrt{10^{4}\times 10^{-4}/400}$ cm
$= \qty{0.05}{cm} = \qty{0.5}{mm}$. Na \qty{2}{mm}: $e^{-4} \approx
2\,\%$. Verdubbeling van $\lambda$ vergt viermaal de diameter,
\qty{4}{\micro m}.
\end{solution}

\begin{solution}{exo:b3:nervous-systems:6}
Sensorisch: $12\times 6 = \qty{72}{m/s}$, $1.2/72 = \qty{17}{ms}$.
Motorisch: \qty{90}{m/s}, $1.1/90 = \qty{12}{ms}$. Twee synapsen,
\qty{1}{ms}: ongeveer \qty{30}{ms}. Ongemyeliniseerd:
$1.1\sqrt{12} = \qty{3.8}{m/s}$ en
$1.1\sqrt{15} = \qty{4.3}{m/s}$: $315 + 258 + 1 \approx \qty{570}{ms}$
--- te traag om de voet te redden.
\end{solution}

\begin{solution}{exo:b3:nervous-systems:7}
$2.4\times 10^{20}/8.6\times 10^{10} \approx 2.8\times 10^{9}$ ATP per
neuron per seconde; bij $2\times 10^{9}$ per puls is dat gemiddeld
ongeveer \qty{1.4}{Hz}. De code moet spaarzaam zijn: weinig pulsen, met
de informatie in welke neuronen vuren en wanneer, niet in hoge tempo's
overal.
\end{solution}

\begin{solution}{exo:b3:nervous-systems:8}
Twee neuronen die elkaar remmen zonder moe te worden vormen een
schakelaar: het eerste dat vuurt onderdrukt het andere zolang de
aandrijving duurt, en er wordt nooit overgedragen. Voor oscillatie is
een traag proces nodig dat de greep van de winnaar verzwakt ---
aanpassing van zijn vuren, verzwakking van zijn remmende synaps, of een
terugstoot van de verliezer na de bevrijding --- en de tijdconstante
daarvan bepaalt de periode.
\end{solution}

\begin{solution}{exo:b3:nervous-systems:9}
L-dopa wordt door de transporter voor grote neutrale aminozuren over de
barrière gedragen; dopamine, polair en geladen, niet. Patiënten nemen
L-dopa (met een remmer van het perifere enzym, zodat het niet vóór de
overtocht wordt omgezet), en het eigen enzym van het brein maakt er
dopamine uit waar die nodig is. Antilichamen, van \qty{150}{kDa}, komen
noch door de nauwe verbindingen noch via een transporter; om ze af te
leveren is een pendel nodig die een receptor voor transcytose kaapt, of
een injectie in het \omterm{prop:b3:nervous-systems:barrier-budget}{hersenvocht}.
\end{solution}

\begin{solution}{exo:b3:nervous-systems:10}
De membraanstroom per lengte-eenheid is de resistieve $V/r_{m}$ plus
de capacitieve $c_{m}\,\partial V/\partial t$; behoud van lading geeft
$\partial I/\partial x = -(V/r_{m} + c_{m}\,\partial V/\partial t)$, en
met $\partial V/\partial x = -r_{i}I$, $\frac{1}{r_{i}}\,\partial^{2}V/
\partial x^{2} = V/r_{m} + c_{m}\,\partial V/\partial t$; vermenigvuldigd met
$r_{m}$: $\lambda^{2}\,\partial^{2}V/\partial x^{2} = V + \tau\,\partial
V/\partial t$. Met $\partial V/\partial t = 0$ keert de vergelijking
van de stelling terug. $\lambda$ is een lengte en $\tau$ een tijd, dus
is $\lambda/\tau$ een snelheid --- de natuurlijke schaal voor hoe snel
een verstoring zich uitbreidt.
\end{solution}

\begin{solution}{exo:b3:nervous-systems:11}
$v \propto \sqrt{d}$: viermaal de snelheid vergt zestienmaal de
diameter, \qty{8}{mm}. Oppervlak $\pi(4\,\text{mm})^{2} \approx \qty{50}{mm^2}$
tegen $\pi\,(10\,\unit{\micro m})^{2} = \qty{314}{\micro m^2}$: een verhouding van
$1.6\times 10^{5}$. Een zenuw van snelle ongemyeliniseerde vezels is
onmogelijk --- één zo'n vezel zou zo dik zijn als de zenuw --- dus een
dier met veel snelle kanalen heeft \omterm{def:b3:nervous-systems:cells}{myeline} nodig, en daarom hebben de
gewervelden (en enkele ongewervelden onafhankelijk) het ontwikkeld.
\end{solution}

\begin{solution}{exo:b3:nervous-systems:12}
$a = 0.4/30^{0.75} = 0.4/12.8 = 0.031$; voor \qty{70000}{g}: $0.031\times
\num{70000}^{0.75} = 0.031\times 4300 \approx \qty{134}{g}$. De
menselijke \qty{1350}{g} is tien maal de voorspelling. De hersenmassa
volgt de lichaamsmassa omdat een groter lichaam meer sensorische en
motorische neuronen nodig heeft, en omdat de neuronen zelf met de
hersengrootte meegroeien; wat de cognitie volgt is het aantal neuronen
in de schors, dat afhangt van hoe dicht zij zijn opeengepakt ---
primaten pakken er meer per gram dan andere zoogdieren, en de mens het
meest van alle.
\end{solution}

\begin{solution}{pb:b3:nervous-systems:1}
\textbf{1.} $\lambda = \sqrt{2\times 10^{4}\,d/400}$: voor $d =
\qty{1e-4}{cm}$ is dat $\sqrt{5\times 10^{-3}} = \qty{0.071}{cm} =
\qty{0.71}{mm}$; voor \qty{4}{\micro m}, \qty{1.41}{mm}. $\tau = 2\times
10^{4}\times 10^{-6} = \qty{20}{ms}$.
\textbf{2.} $5\,e^{-0.3/0.71} = \qty{3.3}{mV}$; $5\,e^{-0.3/1.41} =
\qty{4.0}{mV}$.
\textbf{3.} Verder gelegen invoeren komen verzwakt aan, dus bepaalt de
plaats waar een synaps landt haar gewicht; een potentiaal vervalt in
ongeveer $\tau$, dus tellen twee invoeren alleen op als zij binnen zo'n
\qty{20}{ms} van elkaar vallen --- het neuron is een
gelijktijdigheidsmelder met een venster van \qty{20}{ms}.
\textbf{4.} Gemyeliniseerd \qty{60}{m/s}; ongemyeliniseerd $1.1\sqrt{10} =
\qty{3.5}{m/s}$; om ongemyeliniseerd \qty{60}{m/s} te halen, $d = (60/1.1)^{2}
\approx \qty{3000}{\micro m}$: drie millimeter.
\textbf{5.} $\pi(5)^{2} = \qty{79}{\micro m^2}$ tegen $\pi(1500)^{2} =
7\times 10^{6}\,\unit{\micro m^2}$: ongeveer \num{90000} gemyeliniseerde
axonen in de ruimte van één.
\textbf{6.} \qty{2}{cm} bij \qty{60}{m/s} is \qty{0.3}{ms}; bij
\qty{3.5}{m/s} \qty{5.7}{ms}: zo'n \qty{5}{ms} extra vertraging. Erger
nog, het kale membraan heeft een grote capaciteit en weinig kanalen,
dus kan de stroom van de laatste gave knoop het niet tot de drempel
brengen: geleidingsblok.
\textbf{7.} Sensorisch $1.2/60 = \qty{20}{ms}$; twee synapsen
\qty{1}{ms}; motorisch $1.1/90 = \qty{12}{ms}$: ongeveer \qty{33}{ms}
(plus een milliseconde bij de neuromusculaire overgang).
\textbf{8.} Naar het brein: $0.6/60 = \qty{10}{ms}$ plus twee synapsen,
\qty{11}{ms} nadat het signaal bij \qty{20}{ms} het merg binnenkomt: het
signaal bereikt de schors ongeveer wanneer de voet beweegt; de bewuste
ervaring van pijn vergt nog enkele honderden milliseconden verwerking,
zodat de voet wordt teruggetrokken voordat de pijn wordt gevoeld.
\textbf{9.} $1.1\sqrt{1} = \qty{1.1}{m/s}$: $1.8/1.1 \approx
\qty{1.6}{s}$. De scherpe eerste pijn komt binnen tientallen
milliseconden aan via dunne gemyeliniseerde vezels; de doffe, brandende
tweede pijn een seconde of meer later via de ongemyeliniseerde
C-vezels.
\textbf{10.} Sensorisch $3.2/60 = \qty{53}{ms}$, motorisch $3.1/90 =
\qty{34}{ms}$, synapsen \qty{1}{ms}: ongeveer \qty{90}{ms}. Lange
dieren hebben trage reflexen; zij maken dat goed met dikkere axonen en
door meer aan plaatselijke schakelingen over te laten.
\textbf{11.} $0.8/80 = \qty{10}{ms}$ per richting, \qty{20}{ms}; synaps
\qty{0.5}{ms}, overgang \qty{1}{ms}, kracht \qty{5}{ms}: \qty{26.5}{ms},
waarbij de omzetting in de spoel zelf en de geleiding de spier in de
rest van de \qty{30}{ms} vullen.
\textbf{12.} Ongemyeliniseerde axonen geleiden met een meter of twee per
seconde, dus lussen die bij een volwassene \qty{30}{ms} kosten duren
honderden, met een slechte tijdsordening; de myelinisatie verhoogt de
snelheid tien- tot vijftigvoudig en maakt de tijdsordening nauwkeurig,
en dat is wat lopen, grijpen en spreken vereisen.
\textbf{13.} $20/50\,000 = 4\times 10^{-4}$ mol/s $= 2.4\times 10^{20}$
ATP per seconde; $2.8\times 10^{9}$ per neuron per seconde.
\textbf{14.} $4\times 10^{8} + 7000\times 2\times 10^{4} = 4\times 10^{8}
+ 1.4\times 10^{8} = 5.4\times 10^{8}$ ATP.
\textbf{15.} Alles aan pulsen: $2.8\times 10^{9}/5.4\times 10^{8} \approx
\qty{5}{Hz}$; de helft: ongeveer \qty{2.6}{Hz}.
\textbf{16.} Bij \qty{1}{Hz} is dat $5.4\times 10^{8}/2.8\times 10^{9} \approx
\qty{20}{\%}$ van de begroting: de code is spaarzaam --- de meeste
neuronen zwijgen het grootste deel van de tijd, en de informatie zit in
welke enkele vuren en wanneer.
\textbf{17.} $8.6\times 10^{10}\times 50\times 5.4\times 10^{8} = 2.3
\times 10^{21}$ ATP/s $= 3.9\times 10^{-3}$ mol/s $\approx \qty{190}{W}$
--- tweemaal het hele lichaam in rust.
\textbf{18.} $5\times 16 = \qty{80}{kJ/h} = \qty{22}{W}$: net genoeg,
zonder reserve. Stopt de toevoer, dan stoppen de pompen binnen seconden,
lopen de gradiënten leeg en gaat het bewustzijn in ongeveer tien
seconden verloren.
\textbf{19.} $a = 0.4/12.8 = 0.031$; \qty{70}{kg}: $0.031\times 4300 =
\qty{134}{g}$; \qty{5000}{kg}: $0.031\times 1.06\times 10^{5} \approx
\qty{3300}{g}$.
\textbf{20.} Mens $1350/134 \approx 10$; olifant $4800/3300 \approx
1.5$.
\textbf{21.} In de \omterm{def:b3:nervous-systems:plan}{kleine hersenen} --- \qty{98}{\%} ervan --- waar zij
de veertigduizend spieren van de slurf coördineren. Het aantal neuronen
in de schors volgt de cognitie, niet het totaal; de \omterm{def:b3:nervous-systems:plan}{kleine hersenen}
volgen de motorische taak van het besturen van een enorm lichaam.
\textbf{22.} Axonen en dendrieten worden langer met het brein, dus neemt
elk neuron meer volume in en groeit het aantal neuronen trager dan de
massa. Bij primaten blijft de neurongrootte vrijwel gelijk terwijl het
brein groeit, zodat het aantal neuronen bijna evenredig met de massa
groeit, en een primatenbrein van een gegeven gewicht enkele malen
zoveel neuronen bevat als een knaagdierbrein van dat gewicht.
\textbf{23.} Een arm kan met plaatselijke schakelingen zijn eigen
zuignappen en geledingen coördineren, reiken, grijpen, voedsel langs
zichzelf doorgeven en vermijden dat hij zijn eigen huid grijpt; hij kan
niet zien, geen doel kiezen en geen onderscheid leren. Toets: snijd de
zenuwstreng van de arm van het brein los en raak de arm met voedsel aan
--- hij grijpt het en geeft het door naar waar de mond zou zijn; een
afgezette arm doet hetzelfde nog een uur lang.
\textbf{24.} Ongeveer $23$ synapsen per neuron tegen $7000$: driehonderd
maal minder. De volledige kaart vertelt welk neuron welk ander kan
beïnvloeden, maar niet de sterkte of het teken van elke verbinding,
welke neuromodulatoren ze veranderen, of de dynamiek --- zodat het
bedradingsschema zelfs voor de worm het gedrag maar ten dele voorspelt.
\textbf{25.} $\lambda = \qty{0.71}{mm}$ voor de dendriet van
\qty{1}{\micro m}; latentie van het terugtrekken ongeveer \qty{33}{ms};
gemiddeld vuurtempo ongeveer \qty{2.6}{Hz} met de helft van de
begroting aan pulsen.
\end{solution}
